\documentclass[noauthor,noinstructornotes]{ximera}

\title{Inverse Game Revisited}
\author{Kenneth Berglund}

\begin{document}
\begin{abstract}
    Playing a game to get an even better sense of inverse functions.
\end{abstract}
\maketitle

Let's play a similar game to the one we played two weeks ago! 
You'll pair up to think about this problem together, but play against me this time.
We'll start with a function, 
\(f\left(x\right)=2^x\). I'll select an \(x\)-value in secret and 
tell you the corresponding output. Your job is to guess what my 
\(x\)-value was. 

\begin{enumerate}
    \item If the output is 4, what is my \(x\)-value?
    \item What if my output is 2?
    \item 1?
    \item \(\frac{1}{2}\)?
    \item 64?
    \item \(\sqrt{2}\)?
\end{enumerate}

What process did you use to find these \(x\)-values? Be as specific as possible.

If you get bored, try to play my (the author's) role.

\begin{instructorNotes}
Emphasize that we're drawing on past experience here. This is also a 
warm-up for finding common logarithm values.
\end{instructorNotes}

\end{document}
